Ge\+N-\/\+Foam is launched like any Open\+F\+O\+AM solver, by executing the following commands in a terminal (after sourcing the Open\+F\+O\+AM environment)\+:

{\ttfamily Ge\+N-\/\+Foam}

or

{\ttfamily mpirun -\/np (number of processors) Ge\+N-\/\+Foam -\/parallel}

In case of parallel calculations, one should decompose each one of the three meshes using the command

{\ttfamily decompose\+Par -\/region (region name)}

where the region name is fluid\+Region, neutro\+Region or thermal\+Mechanical\+Region. {\ttfamily decompose\+Par} operates according to the standard Open\+F\+O\+AM {\itshape decompose\+Par\+Dict} to be placed in {\itshape system/(region name)}

N.\+B.\+: all meshes must be decomposed in the same number of domains, and a consistent decompose\+Par\+Dict must also be present in {\itshape system/}

In the tutorials, {\itshape Allrun} (and sometimes {\itshape Allrun\+\_\+parallel}) bash scripts are provided that can be used to run the case (one should just type {\ttfamily Allrun} or {\ttfamily Allrun\+\_\+parallel} in a terminal ), including mesh decomposition for parallel cases, as well as multiple Ge\+N-\/\+Foam runs that are used for instance to\+: 1) achieve a steady-\/state; 2) run a transient starting from that steady-\/state. 